\chapter{模型求解方法}
\label{chap:model-solving}

\begin{quote}\small
\textit{本章摘要。}给定数据、任务和模型以后，还需要计算模型中的未知量。对于确定性参数模型，求解通常得到一个参数值，可以使用闭式计算、梯度方法或启发式搜索。对于随机参数模型，求解需要计算参数的后验分布。本章还讨论含隐变量的极大似然估计：此时参数仍是固定未知量，但参数更新需要隐变量的条件分布。期望最大化算法和变分推断分别说明怎样利用条件分布求参数，以及怎样用可计算的分布近似后验。
\end{quote}

上一章区分了训练目标与评价指标。本章讨论如何依据训练目标计算模型中的未知量。最小二乘回归可以转化为线性方程组。含隐变量的极大似然估计需要计算隐变量的条件分布。贝叶斯学习则需要计算参数后验。求解是这些计算任务的统称，优化是其中一类方法。

设训练数据为$\mathcal D=\{(\bm x_i,y_i)\}_{i=1}^n$。确定性参数模型把$\theta$视为固定未知量，常通过最小化训练目标得到点估计
\begin{equation}
 \hat\theta=\argmin_\theta J(\theta).
 \label{eq:solving-deterministic-objective}
\end{equation}
这里$J$由训练损失和可能的正则项构成。贝叶斯模型为参数指定先验分布，并用$p(\theta\mid\mathcal D)$表示观测数据后的参数后验。具有随机隐变量的模型也可以采用固定参数估计。本章所说的随机变量来自模型的概率设定，与算法是否随机抽取小批量无关。

隐变量$\bm z$是没有被观测、但参与数据生成的随机量。含隐变量的模型既可以用固定参数进行极大似然估计，也可以进行贝叶斯学习。因此，本章先讨论固定参数的直接计算与搜索，再讨论隐变量怎样影响极大似然求解，最后讨论随机参数的后验近似。

\section{确定性参数模型：从目标到参数值}

\subsection{先判断是否存在直接解}
第三章的线性回归采用平方损失：
\begin{equation}
 J(\bm w)=\frac12\|X\bm w-\bm y\|_2^2,
 \qquad \nabla J(\bm w)=X^{\mathsf T}(X\bm w-\bm y).
 \label{eq:solving-least-squares}
\end{equation}
其中$X\in\R^{n\times d}$是设计矩阵，$\bm y\in\R^n$是标签向量，$\bm w\in\R^d$是待定参数。令梯度为零，得到正规方程
\begin{equation}
 X^{\mathsf T}X\hat{\bm w}=X^{\mathsf T}\bm y.
 \label{eq:solving-normal-equation}
\end{equation}
由于平方损失是凸函数，正规方程的解是全局最小点。当$X$列满秩时，解唯一，可以写为$\hat{\bm w}=(X^{\mathsf T}X)^{-1}X^{\mathsf T}\bm y$。

这个表达式说明解的数学形式。实际计算通常对$X$进行QR分解或奇异值分解（Singular Value Decomposition，SVD），直接求解最小二乘问题，避免显式求逆和形成$X^{\mathsf T}X$造成的数值误差放大。列线性相关时解不唯一，可用SVD求最小范数解；列接近相关时，即使解唯一，也可能对数据扰动敏感。正则化改变目标，可以改善稳定性，但其强度仍需验证。

许多非线性模型没有可用的闭式解。有闭式表达式的问题在规模很大时，也可能更适合迭代计算。接下来讨论怎样从一个初始参数出发，逐步降低目标。

\subsection{梯度怎样给出局部改进方向}
若$J$可微，小位移$\Delta$引起的目标变化近似为
\begin{equation}
 J(\theta+\Delta)-J(\theta)\approx\nabla J(\theta)^{\mathsf T}\Delta.
\end{equation}
在欧氏长度固定的位移中，与梯度反向的位移使这个一阶变化最小。因此，梯度下降采用
\begin{equation}
 \theta_{t+1}=\theta_t-\eta_t\nabla J(\theta_t).
 \label{eq:solving-gradient-descent}
\end{equation}
这里$t$是更新次数，$\eta_t>0$是学习率。梯度给出当前位置的局部信息，学习率决定沿该方向走多远。

对$J(\theta)=\theta^2/2$，更新为$\theta_{t+1}=(1-\eta_t)\theta_t$。固定学习率满足$0<\eta<2$时，非零参数的绝对值逐步减小；$\eta=2$时参数在正负之间交替，损失不变；$\eta>2$时，从非零初值出发会发散。选择下降方向以后，还需要控制步长。

更一般地，假设$J$的梯度满足Lipschitz连续条件：存在$L>0$，使
\begin{equation}
 \|\nabla J(\bm u)-\nabla J(\bm v)\|_2\leq L\|\bm u-\bm v\|_2.
\end{equation}
常数$L$限制梯度随位置变化的速度。令$\Delta=\bm v-\bm u$，沿线段积分可得
\begin{align}
 J(\bm u+\Delta)-J(\bm u)
 &=\int_0^1\nabla J(\bm u+s\Delta)^{\mathsf T}\Delta\,ds\\
 &\leq\nabla J(\bm u)^{\mathsf T}\Delta+\frac L2\|\Delta\|_2^2.
\end{align}
第二步把梯度拆成$\nabla J(\bm u)$与梯度差，再用柯西不等式将梯度差的内积上界为$Ls\|\Delta\|_2^2$。代入$\Delta=-\eta\nabla J(\bm u)$，得到
\begin{equation}
 J(\theta_t-\eta\nabla J(\theta_t))
 \leq J(\theta_t)-\eta\left(1-\frac{L\eta}{2}\right)\|\nabla J(\theta_t)\|_2^2.
 \label{eq:solving-descent-bound}
\end{equation}
当$0<\eta<2/L$且$\nabla J(\theta_t)\ne\bm0$时，上式右端减去的量严格为正。因此，更新后的训练目标值严格小于更新前的目标值，即$J(\theta_{t+1})<J(\theta_t)$。这个结论依赖梯度的光滑性和步长范围，只保证当前更新降低目标。对于非凸目标，它不保证达到全局最优，梯度为零的点也可能是局部极小点或鞍点。

实际计算中，可以从一个试探步长开始；如果更新没有满足预定的下降条件，就缩小步长后重新计算，这称为回溯线搜索。达到步数上限、目标变化很小或梯度很小时可以停止，但需要记录具体条件。停止只表示触发了计算规则，不能单独证明已经找到最优解。

\subsection{链式法则与神经网络的梯度计算}

梯度更新要求知道每个参数对目标函数的影响。在线性回归中，这种关系可以直接写出。神经网络的参数先影响某一层的输出，再通过后续各层影响预测值，最后影响损失。要计算较早层的参数梯度，需要把这些连续的影响连接起来。链式法则提供了计算复合函数导数的方法，反向传播则按网络的依赖顺序组织这些计算。

\subsubsection{从一条计算链理解链式法则}
设一个参数$w$先决定中间量$a=g(w)$，再由$a$决定损失$\ell=h(a)$。当两个函数在相应位置可微时，链式法则为
\begin{equation}
 \frac{d\ell}{dw}=\frac{d\ell}{da}\frac{da}{dw}.
 \label{eq:solving-chain-rule}
\end{equation}
其中$da/dw$描述参数变化怎样影响中间量，$d\ell/da$描述中间量变化怎样影响损失。对足够小的变化$\Delta w$，先有$\Delta a\approx(da/dw)\Delta w$，再有$\Delta\ell\approx(d\ell/da)\Delta a$。将两步连接，便得到式\eqref{eq:solving-chain-rule}中的导数乘积。

例如，固定输入$x$和标签$y$，令$a=wx+b$、$\hat y=\phi(a)$、$\ell=(\hat y-y)^2/2$。函数$\phi$称为激活函数，用来对加权和进行非线性变换。沿$w\to a\to\hat y\to\ell$这条计算链，有
\begin{equation}
 \frac{\partial\ell}{\partial w}
 =\frac{\partial\ell}{\partial\hat y}
  \frac{\partial\hat y}{\partial a}
  \frac{\partial a}{\partial w}
 =(\hat y-y)\phi'(a)x.
\end{equation}
三个因子分别来自损失函数、激活函数和加权和。偏置$b$的局部导数为$\partial a/\partial b=1$，所以$\partial\ell/\partial b=(\hat y-y)\phi'(a)$。

网络中一个量可能影响多个后续量。若$\ell=h(u,v)$，而$u$和$v$都依赖$w$，则
\begin{equation}
 \frac{d\ell}{dw}
 =\frac{\partial\ell}{\partial u}\frac{du}{dw}
 +\frac{\partial\ell}{\partial v}\frac{dv}{dw}.
\end{equation}
因此，沿同一条路径的导数相乘，不同路径对同一变量的贡献相加。多层网络和参数共享都需要遵循这条规则。

\subsubsection{一个两层网络的完整计算}
取单个样本$\bm x\in\R^d$和标量标签$y$。考虑一个隐藏层和一个线性输出层：
\begin{align}
 \bm a&=W_1\bm x+\bm b_1,\qquad \bm h=\phi(\bm a),\\
 \hat y&=\bm w_2^{\mathsf T}\bm h+b_2,\qquad
 \ell=\frac12(\hat y-y)^2.
 \label{eq:solving-two-layer-forward}
\end{align}
这里$W_1\in\R^{m\times d}$是第一层权重矩阵，$\bm b_1\in\R^m$是偏置，$\bm a\in\R^m$是加权和，$\bm h\in\R^m$是隐藏层输出。激活函数$\phi$逐坐标作用，$\bm w_2\in\R^m$和$b_2\in\R$是输出层参数。将这些参数放在一起就得到前文的$\theta$。按式\eqref{eq:solving-two-layer-forward}从输入计算到损失，称为前向传播。

求导从损失开始。先计算损失对预测值的导数：
\begin{equation}
 \delta_2=\frac{\partial\ell}{\partial\hat y}=\hat y-y.
\end{equation}
$\delta_2$表示预测值发生微小变化时，损失的一阶变化率。由于$\partial\hat y/\partial w_{2j}=h_j$和$\partial\hat y/\partial b_2=1$，输出层梯度为
\begin{equation}
 \nabla_{\bm w_2}\ell=\delta_2\bm h,\qquad
 \frac{\partial\ell}{\partial b_2}=\delta_2.
\end{equation}
再把导数传到隐藏层。第$j$个隐藏单元满足
\begin{equation}
 \frac{\partial\ell}{\partial h_j}=\delta_2w_{2j},\qquad
 \frac{\partial\ell}{\partial a_j}=\delta_2w_{2j}\phi'(a_j).
\end{equation}
令$\bm\delta_1=\nabla_{\bm a}\ell$，把各坐标放在一起，得到
\begin{equation}
 \bm\delta_1=(\bm w_2\delta_2)\odot\phi'(\bm a).
\end{equation}
符号$\odot$表示逐坐标相乘。最后，由$a_j=\sum_{k=1}^d W_{1,jk}x_k+b_{1j}$，有$\partial a_j/\partial W_{1,jk}=x_k$，于是
\begin{equation}
 \frac{\partial\ell}{\partial W_{1,jk}}=\delta_{1j}x_k,\qquad
 \nabla_{W_1}\ell=\bm\delta_1\bm x^{\mathsf T},\qquad
 \nabla_{\bm b_1}\ell=\bm\delta_1.
 \label{eq:solving-two-layer-backward}
\end{equation}
外积$\bm\delta_1\bm x^{\mathsf T}$的形状为$m\times d$，与$W_1$一致。其余梯度也与对应参数形状一致。计算从损失向输入方向进行，逐层复用已经得到的导数，这就是反向传播（Backpropagation）。

取一个输入和一个隐藏单元，令$x=2$、$y=1$、$W_1=w_2=1$、$b_1=b_2=0$，使用修正线性单元（Rectified Linear Unit，ReLU）$\phi(a)=\max(0,a)$。前向计算得到$a=2$、$h=2$、$\hat y=2$和$\ell=0.5$。由于$a>0$，$\phi'(a)=1$。反向计算得到$\delta_2=1$、$\delta_1=1$，两个权重的梯度均为$2$，两个偏置的梯度均为$1$。

若采用学习率$0.1$同时更新所有参数，则$W_1=w_2=0.8$、$b_1=b_2=-0.1$。重新前向计算得到$\hat y=0.8(0.8\times2-0.1)-0.1=1.1$，损失为$0.005$。这个例子展示了从前向计算、反向求导到参数更新的完整过程；它不意味着任意网络和步长都能降低损失。

\subsubsection{多层网络怎样复用导数}
对具有$L$层的网络，令$\bm h^{(0)}=\bm x$，第$l$层计算
\begin{equation}
 \bm a^{(l)}=W^{(l)}\bm h^{(l-1)}+\bm b^{(l)},\qquad
 \bm h^{(l)}=\phi_l(\bm a^{(l)}).
\end{equation}
若第$l$层宽度为$d_l$，则$W^{(l)}\in\R^{d_l\times d_{l-1}}$，$\bm a^{(l)}$、$\bm h^{(l)}$和$\bm b^{(l)}$均属于$\R^{d_l}$。这里假设激活函数逐坐标作用，最终损失为$\ell(\bm h^{(L)},y)$。定义$\bm\delta^{(l)}=\nabla_{\bm a^{(l)}}\ell$，输出层先计算
\begin{equation}
 \bm\delta^{(L)}=\nabla_{\bm h^{(L)}}\ell\odot\phi_L'(\bm a^{(L)}).
\end{equation}
随后对$l=L-1,\ldots,1$依次计算
\begin{align}
 \bm\delta^{(l)}
 &=\left((W^{(l+1)})^{\mathsf T}\bm\delta^{(l+1)}\right)
   \odot\phi_l'(\bm a^{(l)}),\\
 \nabla_{W^{(l)}}\ell
 &=\bm\delta^{(l)}(\bm h^{(l-1)})^{\mathsf T},\qquad
 \nabla_{\bm b^{(l)}}\ell=\bm\delta^{(l)}.
 \label{eq:solving-backprop-recursion}
\end{align}
输出层的参数梯度也使用第二行公式。矩阵转置乘法把下一层各单元的导数贡献加到当前层，再乘当前激活的导数。对于跨坐标耦合的运算，需要使用相应的矩阵导数，不能直接套用逐坐标激活公式。

前向传播需要保存求导所需的输入和中间结果。反向传播按依赖的逆序计算局部导数，并累加分支贡献。这样无需为每个参数重新计算整条导数链。自动微分（Automatic Differentiation，AD）系统根据实际执行的运算记录这些依赖，并按局部求导规则完成累积。神经网络训练通常使用反向模式自动微分，它适合从一个标量损失计算大量参数的梯度。

上述推导在相应运算可微处成立。ReLU在零点不可微，实际实现通常约定该点的局部导数为零；包含这类运算时，前文的光滑下降保证不能未经检查直接应用。对训练集平均损失，逐样本梯度也取平均，正则项的梯度另行加入。所有梯度必须在同一组更新前参数下算完，再由优化器更新参数。反向传播解决梯度怎样计算的问题，梯度下降、动量或Adam解决怎样使用梯度更新参数的问题。

\subsection{随机梯度、动量与自适应尺度}
设$J(\theta)=n^{-1}\sum_{i=1}^n\ell_i(\theta)$，其中$\ell_i$是第$i$个样本的损失。大数据集上每步计算完整梯度成本较高。随机梯度下降（Stochastic Gradient Descent，SGD）均匀抽取非空小批量$B_t$，计算
\begin{equation}
 \hat{\bm g}_t=\frac{1}{|B_t|}\sum_{i\in B_t}\nabla_\theta\ell_i(\theta_t),
 \qquad \theta_{t+1}=\theta_t-\eta_t\hat{\bm g}_t.
 \label{eq:solving-sgd}
\end{equation}
在固定当前参数、按上述方式抽样时，$\E[\hat{\bm g}_t\mid\theta_t]=\nabla J(\theta_t)$。小批量梯度有噪声，单步目标可能增加，但每步所需样本较少。若目标还含正则项，应加上其梯度。把批平均改为批求和会改变更新尺度，必须相应调整学习率。

如果参数在某些方向来回移动，可以对梯度进行指数加权平均。令$\bm v_0=\bm0$，从$t=1$开始更新：
\begin{align}
 \bm v_t&=\beta\bm v_{t-1}+(1-\beta)\hat{\bm g}_t,\\
 \theta_{t+1}&=\theta_t-\eta_t\bm v_t.
\end{align}
这里$0\leq\beta<1$，$\hat{\bm g}_t$在$\theta_t$处计算。持续同向的梯度会被保留，正负交替的梯度会部分抵消。这是动量方法的一种约定。有些实现省略$(1-\beta)$，对应的有效步长不同。

不同坐标的梯度尺度也可能差异很大。自适应矩估计（Adaptive Moment Estimation，Adam）同时维护梯度与梯度平方的指数平均\cite{kingma2015adam}：
\begin{align}
 \bm m_t&=\beta_1\bm m_{t-1}+(1-\beta_1)\hat{\bm g}_t,\\
 \bm s_t&=\beta_2\bm s_{t-1}+(1-\beta_2)\hat{\bm g}_t\odot\hat{\bm g}_t,\\
 \hat{\bm m}_t&=\bm m_t/(1-\beta_1^t),\qquad
 \hat{\bm s}_t=\bm s_t/(1-\beta_2^t),\\
 \theta_{t+1}&=\theta_t-\eta_t\frac{\hat{\bm m}_t}{\sqrt{\hat{\bm s}_t}+\epsilon}.
 \label{eq:solving-adam}
\end{align}
初始状态为$\bm m_0=\bm s_0=\bm0$，$0\leq\beta_1,\beta_2<1$，$\epsilon>0$用于数值稳定。符号$\odot$表示逐坐标相乘，除法和平方根也逐坐标进行。

零初始化使早期平均偏小。例如，恒定梯度$\bm g$对应$\bm m_t=(1-\beta_1^t)\bm g$，除以$1-\beta_1^t$便恢复原尺度。平方梯度的校正同理。归一化使近期平方梯度较大的坐标使用较小的更新尺度。这里的二阶矩是梯度平方的平均，与描述目标曲率的海森矩阵（Hessian matrix）不同。Adam仍需要选择学习率和停止规则，其更新不保证每一步目标下降。

\subsection{启发式搜索求解模型参数}
离散规则、结构选择和仿真程序可能无法提供可用梯度。此时可以直接评价候选参数的目标值，再按搜索规则产生新候选。启发式优化适合这类问题，其效果取决于候选表示、搜索规则和计算预算。

下面统一考虑最小化$J(\theta)$，并要求候选属于可行集合$\Omega$。算法只需要能够评价候选的目标值。遗传算法通过组合与变异产生新候选，模拟退火通过邻域移动搜索，粒子群则利用多个候选已经找到的较好位置指导移动。

\paragraph{遗传算法。}
遗传算法（Genetic Algorithm，GA）同时维护$M$个候选，组成第$t$代种群$\{\theta_t^{(i)}\}_{i=1}^M$。每个候选也称为个体。算法首先选择参数的编码方式：连续参数可以直接编码为实数向量，特征是否被选用可以编码为二进制向量。编码决定交叉和变异怎样进行，因此需要先明确哪些编码对应合法模型。

每代先计算$J_i=J(\theta_t^{(i)})$，再根据目标值选择父代。选择规则应使较优个体更可能被选中，同时保留其他个体参与搜索的机会。例如，可以将目标值转换为正的选择权重：
\begin{equation}
 s_i=\exp\left(-\frac{J_i-J_{\min}}{\tau}\right),\qquad
 P_i=\frac{s_i}{\sum_{j=1}^M s_j},\qquad J_{\min}=\min_j J_j.
 \label{eq:solving-ga-selection}
\end{equation}
这里$P_i$是选择第$i$个个体作为父代的概率，$\tau>0$控制选择压力，与目标值使用相同尺度。目标较小的个体具有较大的选择概率。$\tau$较小时，选择更集中于少数较优个体；$\tau$较大时，各个体的概率更接近。这个公式是一种可用选择规则，实际也常采用随机抽取若干个体、选其中最优者的锦标赛选择。

选出两个父代$\theta^{(a)}$和$\theta^{(b)}$后，通过交叉生成子代。对实数编码，一种算术交叉为
\begin{equation}
 \widetilde\theta=\alpha\theta^{(a)}+(1-\alpha)\theta^{(b)},\qquad
 \alpha\sim\operatorname{Uniform}(0,1).
\end{equation}
子代位于两个父代之间。若只进行这种交叉，新候选无法离开已有个体的凸包，因此还需要变异。对$d$维实数参数，可以使用
\begin{equation}
 \theta'=\widetilde\theta+\sigma\bm m\odot\epsilon,\qquad
 m_j\sim\operatorname{Bernoulli}(p_{\mathrm{mut}}),\quad
 \epsilon\sim\mathcal N(\bm0,I).
 \label{eq:solving-ga-mutation}
\end{equation}
这里各坐标的$m_j$独立，$p_{\mathrm{mut}}$是变异概率，$\sigma>0$是扰动尺度。$m_j=0$时保持该坐标，$m_j=1$时加入随机扰动。若各参数单位不同，应分别设置坐标尺度。交叉与变异产生的候选必须检查约束，必要时修复或重新生成。

例如，从四个特征中选择恰好两个，可以用$(1,0,1,0)$表示选择第一和第三个特征。交换一个已选位置和一个未选位置是合法变异。单点交叉把两个父代在某位置切开，再拼接前后片段，但可能产生选择数量错误的子代，需要修复。二进制编码不能直接使用上述实数交叉和高斯变异。

生成并评价足够多的子代后，组成下一代，并可保留上一代的最优个体，这称为精英保留。在目标评价固定的条件下，它使历代最优目标值不增加，但不保证达到全局最优。种群大小决定每代评价成本，选择压力过强会使个体过早相似，变异过大则可能使搜索难以保留已有结构。遗传算法适用于离散结构、组合规则或混合编码问题，需按评价预算或最优值长期无改进等条件停止。

\paragraph{模拟退火。}
模拟退火（Simulated Annealing，SA）维护一个当前候选$\theta_t$，通过局部修改提出$\theta'$。只接受改进候选的搜索容易停在局部极小点附近。模拟退火允许以一定概率接受较差候选，使搜索有机会经过目标较高的区域，再找到其他较低区域。

先指定邻域提议规则。例如，无约束连续参数可以采用$\theta'=\theta_t+\sigma\epsilon$，其中$\epsilon\sim\mathcal N(\bm0,I)$；固定数量的特征选择可以交换一个已选特征和一个未选特征。对于对称提议，即从$\theta_t$提出$\theta'$与反向提出的概率相同，令$\Delta=J(\theta')-J(\theta_t)$，接受概率为
\begin{equation}
 a(\theta_t,\theta')=\min\{1,\exp(-\Delta/T_t)\}.
 \label{eq:solving-simulated-annealing}
\end{equation}
这里$T_t>0$是温度，与目标差值使用相同尺度。每次抽取$u\sim\operatorname{Uniform}(0,1)$，若$u\leq a(\theta_t,\theta')$就令$\theta_{t+1}=\theta'$，否则令$\theta_{t+1}=\theta_t$。当$\Delta\leq0$时总是接受；当$\Delta>0$时，目标增加越多，接受概率越小。

固定温度下，这一接受规则可由希望采样的分布$\pi_T(\theta)\propto\exp(-J(\theta)/T)$理解。当该分布可归一化时，对称提议下两个候选的密度比就是$\exp(-\Delta/T)$。较小温度使概率更集中于目标较小的区域。例如，$\Delta=2$时，$T_t=2$对应接受概率约$0.368$，$T_t=0.5$时约$0.018$。若提议不对称，维持上述固定温度分布还需要乘以反向与正向提议的概率比，不能直接使用式\eqref{eq:solving-simulated-annealing}。

搜索开始时设置初温$T_0$，随后逐渐降温。常用几何降温：
\begin{equation}
 T_k=T_0\gamma^k,\qquad 0<\gamma<1.
\end{equation}
这里$k$表示温度阶段，每个阶段可以进行多次邻域提议。较高温度允许更多非改进移动，较低温度使搜索更偏向改进。初温应结合典型目标差值选择，降温速度应结合评价预算选择。邻域过小可能移动缓慢，过大可能频繁提出很差的候选。

由于算法允许接受较差候选，当前目标值可能上升，应另行保存搜索中找到的最优候选。适当条件和很慢的降温计划下存在渐近全局收敛结论，但常用有限预算的几何降温不提供普遍的全局最优保证。模拟退火适合容易定义合法局部修改的问题，其局限是单个候选可能需要较长时间探索多个区域。

\paragraph{粒子群优化。}
粒子群优化（Particle Swarm Optimization，PSO）维护多个连续参数候选，每个候选称为粒子。与遗传算法通过交叉产生子代不同，粒子群保留各粒子的位置和移动速度，并利用历史较好位置指导下一次移动。它适用于能够直接评价目标、但梯度难以取得的连续参数搜索。

令第$i$个粒子在第$t$轮的位置为$\theta_t^{(i)}\in\R^d$，速度为$\bm v_t^{(i)}\in\R^d$。记录该粒子截至当前找到的最优位置$\bm p_t^{(i)}$，并从所有粒子的历史最优位置中选择群体最优位置$\bm g_t$：
\begin{equation}
 \bm p_t^{(i)}\in\argmin_{\theta\in\{\theta_0^{(i)},\ldots,\theta_t^{(i)}\}}J(\theta),\qquad
 \bm g_t\in\argmin_{\theta\in\{\bm p_t^{(1)},\ldots,\bm p_t^{(M)}\}}J(\theta).
\end{equation}
这里$M$是粒子数量。初始化时先评价各粒子，令$\bm p_0^{(i)}=\theta_0^{(i)}$，再确定$\bm g_0$。一种常见更新规则为
\begin{align}
 \bm v_{t+1}^{(i)}
 &=\omega\bm v_t^{(i)}
 +c_1\bm r_{1,t}^{(i)}\odot(\bm p_t^{(i)}-\theta_t^{(i)})
 +c_2\bm r_{2,t}^{(i)}\odot(\bm g_t-\theta_t^{(i)}),\\
 \theta_{t+1}^{(i)}&=\theta_t^{(i)}+\bm v_{t+1}^{(i)}.
 \label{eq:solving-pso-update}
\end{align}
随机向量$\bm r_{1,t}^{(i)}$和$\bm r_{2,t}^{(i)}$的坐标独立服从$\operatorname{Uniform}(0,1)$，$\omega\geq0$是惯性权重，$c_1,c_2\geq0$是引导强度。第一项保留原来的移动方向，第二项把粒子拉向自身历史最优位置，第三项把粒子拉向群体历史最优位置。两个随机向量使不同粒子对相同位置的移动幅度有所差别。

例如，一维粒子当前位置为$2$、速度为$0.5$，自身历史最优位置为$1$，群体历史最优位置为$0$。若$\omega=0.5$、$c_1=c_2=1$，这次抽到$r_1=r_2=0.5$，则新速度为$0.25-0.5-1=-1.25$，新位置为$0.75$。算法并未计算$J'(2)$，移动方向由已经评价过的较好位置决定。新位置是否更好，必须重新计算$J(0.75)$才能判断。

每轮移动后，评价新位置。若它优于该粒子的历史最优位置，就更新$\bm p^{(i)}$；再根据全部历史记录更新$\bm g$。粒子当前位置的目标可以增加，而保存的群体历史最优目标在固定评价下不会增加。对于箱形约束，可以截断或反射越界位置，并相应处理速度。各参数应使用合理尺度，否则数值范围大的坐标可能主导移动。

惯性或引导强度过大可能造成震荡和越界，所有粒子过早聚集又会降低探索范围。速度上限、局部邻域最优引导或重启部分粒子可以调整搜索行为，但不能普遍保证全局最优。上述规则针对连续参数，离散问题需要专门的编码和更新方式。粒子数量增加会增加每轮目标评价次数，因此应结合维度和评价成本设置预算。

三种方法都通过目标评价比较候选。遗传算法侧重种群的选择与重组，模拟退火侧重单个候选的邻域移动与温度控制，粒子群侧重历史最优位置对连续移动的引导。比较它们时，应统一可行集合、目标评价方式和总评价预算。目标可微且梯度成本可控时，高维连续模型通常更适合梯度方法。

\subsection{含隐变量的极大似然：EM算法}

\subsubsection{为什么完整数据容易求解，观测数据却难求解}
确定性参数的概率模型也可能包含未被观测的随机变量，称为隐变量或潜变量（Latent Variables）。例如，高斯混合模型中，我们观察到了样本，却不知道每个样本来自哪个成分。如果成分归属也被观测到，就可以按成分分别估计均值和方差。缺少这部分信息以后，参数估计与成分归属判断相互依赖：需要知道归属才能方便地估计参数，又需要参数才能判断归属。

神经网络中的隐藏层也没有被直接观测，但在普通前馈网络中，给定输入和参数就能确定隐藏层输出。概率隐变量则允许在同一观测条件下具有多种可能取值，计算观测概率时需要将这些可能性按概率汇总。因此，不能直接用一次前向计算得到的中间值替代隐变量的条件分布。

为统一记号，以下用$\bm x$表示全部观测数据，$\bm z$表示全部隐变量，$\theta$仍是固定未知参数。给定联合模型$p(\bm x,\bm z\mid\theta)$，观测数据对数似然为
\begin{equation}
 \ell(\theta)=\log p(\bm x\mid\theta)
 =\log\int p(\bm x,\bm z\mid\theta)\,d\bm z.
 \label{eq:solving-marginal-likelihood}
\end{equation}
离散隐变量时，积分改为求和。这里$\ell$表示对数似然，与前文的单样本损失按上下文区分。对于条件预测模型，相应关系是$p(y\mid\bm x,\theta)=\int p(y,\bm z\mid\bm x,\theta)\,d\bm z$，隐变量的分布可以依赖输入和参数。

极大似然估计要最大化式\eqref{eq:solving-marginal-likelihood}。困难既可能是积分难以计算，也可能是积分或求和外的对数使最大化变得困难。高斯混合模型的求和可以计算，但对数中的多个成分仍耦合在一起。相比之下，完整数据对数似然$\log p(\bm x,\bm z\mid\theta)$往往能拆成容易处理的项。期望最大化（Expectation-Maximization，EM）算法由此提出一个思路：利用当前参数推断缺失信息，再用这些信息改进参数。下面推导为什么应当对缺失信息取条件期望，以及这种交替计算为什么有效。

\subsubsection{怎样构造一个能逐步提高的下界}
只把隐变量替换为一个最可能取值，会丢掉其他取值的概率，通常不能保证观测似然提高。我们引入一个归一化分布$q(\bm z)$，对所有可能的隐变量取加权平均。暂时不指定$q$，先看它能否把难处理的对数积分变成完整数据对数似然的期望。

对数是凹函数。Jensen不等式说明，对正随机量$U$有$\log\E[U]\geq\E[\log U]$。在$q$覆盖联合分布有关支撑、所需期望有限的条件下，将积分改写为$q$下的期望：
\begin{align}
 \ell(\theta)
 &=\log\int q(\bm z)\frac{p(\bm x,\bm z\mid\theta)}{q(\bm z)}\,d\bm z\\
 &=\log\E_q\left[\frac{p(\bm x,\bm z\mid\theta)}{q(\bm z)}\right]\\
 &\geq\E_q\left[\log\frac{p(\bm x,\bm z\mid\theta)}{q(\bm z)}\right]
 =F(q,\theta).
 \label{eq:solving-em-jensen}
\end{align}
这里
\begin{equation}
 F(q,\theta)=\E_q[\log p(\bm x,\bm z\mid\theta)]+\mathcal H(q),\qquad
 \mathcal H(q)=-\E_q[\log q(\bm z)].
\end{equation}
$\mathcal H(q)$是熵，连续变量时为微分熵。$F$是观测对数似然的下界，也称证据下界（Evidence Lower Bound，ELBO）。固定$q$以后，熵与$\theta$无关，最大化这个下界只需要最大化完整数据对数似然的期望。

但任意下界提高，并不意味着原目标提高。例如，一个下界可能从很低的位置上升，仍低于原目标的旧值。要保证原目标不下降，需要让下界在当前参数$\theta^{(t)}$处与原目标相等，然后提高这个下界。

Jensen不等式在上述比值为常数时取等号，即$p(\bm x,\bm z\mid\theta^{(t)})=Cq(\bm z)$在$q$下几乎处处成立。对两边积分，利用$\int q=1$，得到$C=p(\bm x\mid\theta^{(t)})$。因此，令下界接触当前似然的分布应当是
\begin{equation}
 q_t(\bm z)=\frac{p(\bm x,\bm z\mid\theta^{(t)})}{p(\bm x\mid\theta^{(t)})}
 =p(\bm z\mid\bm x,\theta^{(t)}).
\end{equation}
这就解释了为什么使用当前参数下的隐变量后验：它让构造的下界在当前参数处恰好等于观测对数似然。

\subsubsection{由下界推导E步和M步}
给定初始参数$\theta^{(0)}$，E步先按当前参数计算隐变量后验$q_t$，再形成
\begin{equation}
 Q(\theta\mid\theta^{(t)})=\E_{q_t}[\log p(\bm x,\bm z\mid\theta)].
 \label{eq:solving-em-q}
\end{equation}
期望中的$\theta$是接下来要改变的候选参数；分布$q_t$中的$\theta^{(t)}$是固定的旧参数。计算时必须区分这两者。很多模型只需计算有关充分统计量的条件期望，无需显式保存整个后验分布。

M步固定$q_t$，最大化下界。因为$\mathcal H(q_t)$与候选参数无关，更新为
\begin{equation}
 \theta^{(t+1)}=\argmax_\theta F(q_t,\theta)
 =\argmax_\theta Q(\theta\mid\theta^{(t)}).
 \label{eq:solving-em-m}
\end{equation}
由此得到EM的两个步骤：E步使下界在旧参数处与原目标相等，M步利用这个下界求新参数。得到新参数以后，旧$q_t$通常已不再是新参数下的后验，需要重新执行E步，再继续更新。

从联合目标$F(q,\theta)$看，EM是一种交替最大化：固定$\theta$时选择精确后验作为$q$，固定$q$时更新$\theta$。每一步只处理一部分未知量，因此可以利用完整数据模型中较容易的计算。M步可能有闭式解，也可能需要数值优化。E步则要求后验或有关条件期望能够计算，EM并不能自动解决所有不可计算的积分。

\subsubsection{为什么每轮更新不会降低观测似然}
先用下界说明单调性。E步保证$F(q_t,\theta^{(t)})=\ell(\theta^{(t)})$，M步保证$F(q_t,\theta^{(t+1)})\geq F(q_t,\theta^{(t)})$，而下界始终不超过相应似然。因此
\begin{equation}
 \ell(\theta^{(t+1)})\geq F(q_t,\theta^{(t+1)})
 \geq F(q_t,\theta^{(t)})=\ell(\theta^{(t)}).
\end{equation}
这条不等式连接了旧目标值、提高后的下界和新目标值。它保证观测对数似然不下降；若M步严格提高$Q$，则观测对数似然也严格提高。

还可以直接从条件概率分解看清楚两个变化项。由$p(\bm x,\bm z\mid\theta)=p(\bm x\mid\theta)p(\bm z\mid\bm x,\theta)$，取对数并在固定的$q_t$下求期望，得到
\begin{equation}
 \ell(\theta)=Q(\theta\mid\theta^{(t)})-A(\theta\mid\theta^{(t)}),\qquad
 A(\theta\mid\theta^{(t)})=\E_{q_t}[\log p(\bm z\mid\bm x,\theta)].
\end{equation}
$\ell(\theta)$与$\bm z$无关，左边的期望仍为自身。这里$A$是后验对数密度的期望，只有在$\theta=\theta^{(t)}$时，才等于$-\mathcal H(q_t)$。

定义Kullback--Leibler散度（Kullback--Leibler divergence，KL）为$\KL(q\|p)=\E_q[\log(q/p)]$。在有关期望有定义时，它非负且通常不对称。于是
\begin{align}
 A(\theta^{(t+1)}\mid\theta^{(t)})-A(\theta^{(t)}\mid\theta^{(t)})
 &=\E_{q_t}\left[\log\frac{p(\bm z\mid\bm x,\theta^{(t+1)})}{q_t(\bm z)}\right]\\
 &=-\KL(q_t\|p(\bm z\mid\bm x,\theta^{(t+1)}))\leq0.
\end{align}
M步使$Q$不减，而被减去的$A$不增，因此
\begin{equation}
 \ell(\theta^{(t+1)})-\ell(\theta^{(t)})
 =Q(\theta^{(t+1)}\mid\theta^{(t)})-Q(\theta^{(t)}\mid\theta^{(t)})
 +\KL(q_t\|p(\bm z\mid\bm x,\theta^{(t+1)}))\geq0.
\end{equation}
这个分解说明，新似然的增加来自完整数据期望目标的改进，以及旧隐变量后验与新隐变量后验之间的差异。

对一般$q$，同样的代数还给出下界的精确间隙：
\begin{equation}
 \ell(\theta)=F(q,\theta)+\KL(q\|p(\bm z\mid\bm x,\theta)).
 \label{eq:solving-em-decomposition}
\end{equation}
KL为零时，下界与似然相等。这个恒等式与Jensen推导相互对应，也将引出后面的变分推断。

单调性与收敛是不同结论。如果对数似然有上界，不下降的似然值序列会收敛，但这本身不保证参数序列收敛或达到全局最优。在适当正则条件下，EM的收敛点可为驻点，初始化仍会影响结果。无约束高斯混合模型可能因成分方差趋近零而出现似然无界，不能直接套用似然有界时的结论。

\subsubsection{哪些近似仍保留上述保证}
若E步精确计算后验，而M步只使$Q$不下降，不要求达到其最大值，就得到广义EM（Generalized EM，GEM）。例如，可以在M步使用带下降检查的梯度法最小化$-Q$。上述单调性证明仍成立。

若隐变量后验无法精确计算，则是另一个计算困难。将$q$限制在可处理的分布族内，交替提高$F(q,\theta)$，得到变分EM（Variational EM，VEM）。此时E步通常不能令KL为零，所以提高下界不再自动保证观测似然不下降。用采样估计E步期望则得到蒙特卡洛EM（Monte Carlo EM，MCEM），有限采样误差也可能破坏逐轮单调性。这些方法分别近似分布与期望，不能与只放宽M步的广义EM混为一谈。

\subsection{例子：一维高斯混合模型}
设独立样本$x_1,\ldots,x_n$来自$K$个高斯成分：
\begin{equation}
 p(x_i\mid\theta)=\sum_{k=1}^K\pi_k\mathcal N(x_i\mid\mu_k,\sigma_k^2),
 \qquad \pi_k\geq0,\quad \sum_{k=1}^K\pi_k=1.
\end{equation}
参数$\theta$包含混合权重$\pi_k$、均值$\mu_k$和正方差$\sigma_k^2$。隐变量$z_i=k$表示样本来自第$k$个成分。若成分归属已知，每个成分可以分别拟合；归属未知时，用条件概率作为软计数。

E步由Bayes公式计算责任度：
\begin{equation}
 r_{ik}=P(z_i=k\mid x_i,\theta^{(t)})
 =\frac{\pi_k^{(t)}\mathcal N(x_i\mid\mu_k^{(t)},(\sigma_k^2)^{(t)})}
 {\sum_{j=1}^K\pi_j^{(t)}\mathcal N(x_i\mid\mu_j^{(t)},(\sigma_j^2)^{(t)})}.
\end{equation}
每个样本满足$\sum_k r_{ik}=1$。若两个成分权重相同、方差均为$1$、均值为$0$和$2$，则$x_i=1$的责任度各为$1/2$，而$x_i=0$属于第一成分的责任度为$1/(1+e^{-2})\approx0.881$。

固定责任度，完整数据对数似然的期望为
\begin{equation}
 Q=\sum_{i,k}r_{ik}\left[\log\pi_k-\frac12\log(2\pi\sigma_k^2)
 -\frac{(x_i-\mu_k)^2}{2\sigma_k^2}\right].
\end{equation}
令$N_k=\sum_i r_{ik}$。在$N_k>0$且更新方差为正时，求导并结合权重和为一的约束，得到M步：
\begin{align}
 \pi_k^{(t+1)}&=N_k/n,\\
 \mu_k^{(t+1)}&=N_k^{-1}\sum_i r_{ik}x_i,\\
 (\sigma_k^2)^{(t+1)}&=N_k^{-1}\sum_i r_{ik}(x_i-\mu_k^{(t+1)})^2.
\end{align}
均值公式来自$\sum_i r_{ik}(x_i-\mu_k)=0$。方差公式来自对$\sigma_k^2$求导后得到的$N_k\sigma_k^2=\sum_i r_{ik}(x_i-\mu_k)^2$。权重约束的拉格朗日乘子给出$\pi_k=N_k/n$。它们分别是软样本比例、加权均值和加权方差。

反复执行两步，直到观测对数似然的变化小于约定阈值或达到迭代上限。实践中需要处理空成分与方差退化，例如预先限制方差下界，并在这一约束下完成M步。多个初始值可以得到不同解，应在相同模型和约束下比较似然。成分数的选择和泛化性能还需要独立的模型评价。

\section{随机参数模型的求解}

以上各节主要讨论固定未知参数的求解，结果通常是一个点估计$\hat\theta$。随机参数模型则用概率分布表达对未知参数$\theta$的不确定性。这不要求真实参数在每次观测时重新变化。数据较少时，多个参数取值可能都能解释观测结果；只保留一个参数点，就无法描述这些取值之间的差异及其可信程度。因此，求解需要从参数值扩展到参数分布。

一个模型可以同时包含随机参数和固定未知参数。例如，可以用随机变量表示设备状态，用具有固定未知权重的神经网络表示状态与观测之间的条件分布。前者需要条件分布，后者需要参数估计；两类未知量可以通过交替计算或同一个目标下的联合优化求解。

后验难以精确计算时，变分推断（Variational Inference，VI）选择一个便于计算的分布来近似它，再优化这个分布的参数。这里的优化变量描述的是近似分布，求解结果仍是一个分布。前文的变分EM也使用这种思路近似隐变量后验；蒙特卡洛EM则用采样估计条件期望。下面先写出参数后验，再分别讨论变分近似与后验采样。

\subsection{后验分布：数据怎样更新参数的不确定性}
设观测数据为$\mathcal D$。先验分布$p(\theta)$表达观察当前数据之前对参数取值的判断；似然$p(\mathcal D\mid\theta)$表示给定参数时，模型赋予观测数据的概率或概率密度。观察数据后，贝叶斯公式将二者结合为后验分布（posterior distribution）：
\begin{equation}
 p(\theta\mid\mathcal D)=\frac{p(\mathcal D\mid\theta)p(\theta)}{p(\mathcal D)}.
 \label{eq:solving-posterior}
\end{equation}
后验回答：在已观察到这些数据的条件下，各种参数取值有多可信？分子把先验权重与数据提供的似然相乘，分母$p(\mathcal D)$是归一化常数。它由联合分布对参数积分得到：
\begin{equation}
 p(\mathcal D)=\int p(\mathcal D,\theta)\,d\theta
 =\int p(\mathcal D\mid\theta)p(\theta)\,d\theta.
 \label{eq:solving-posterior-normalizer}
\end{equation}
当参数维度较高，或者似然与先验的形式较复杂时，这个积分通常没有解析形式；即使能够计算它，后验本身以及后验预测中的积分也可能仍然难以直接求出。因此，随机参数模型的核心计算问题，是近似后验分布，或者直接计算关于后验分布的期望。

这产生了两条基本路线。一条是分布近似：选择一个可计算的分布族，用$q_\phi(\theta)$近似$p(\theta\mid\mathcal D)$，通过优化确定变分参数$\phi$。另一条是采样近似：构造以后验为目标分布的采样方法，用参数样本近似后验期望。前者给出一个显式的近似分布，后者未必写出后验的解析表达式，但可以用样本完成预测和其他期望计算。

后验的重要性在于它保留了参数的不确定性。我们可以计算后验均值（或者最大后验估计）作为点估计，也可以计算参数区间的后验概率，说明数据支持的取值范围。预测新标签时，还可以让不同参数按其后验权重共同参与预测：
\begin{equation}
 p(y_*\mid\bm x_*,\mathcal D)=\int p(y_*\mid\bm x_*,\theta)p(\theta\mid\mathcal D)\,d\theta.
 \label{eq:solving-posterior-predictive}
\end{equation}
其中$\bm x_*$是新输入，$y_*$是待预测标签；假设给定输入和参数后，新标签与训练数据条件独立，并将输入作为已知条件。积分把每个参数下的预测按后验密度加权汇总，使参数不确定性进入预测结果。后验较集中时，单点预测可能是合适的近似；后验分散或具有多个峰时，一个参数点可能遗漏其他有较大后验权重的预测。后验是否可信仍取决于先验、似然和数据假设，分布形式本身不保证预测校准。

\subsection{边际似然：模型整体怎样支持观测数据}
式\eqref{eq:solving-posterior}的分母由联合分布对参数积分得到：
\begin{equation}
 p(\mathcal D)=\int p(\mathcal D,\theta)\,d\theta
 =\int p(\mathcal D\mid\theta)p(\theta)\,d\theta.
 \label{eq:solving-model-evidence}
\end{equation}
这个量称为边际似然（marginal likelihood），也称模型证据（model evidence）。这里的“边际”表示通过积分消去参数，只留下观测数据的分布；离散参数时，积分改为求和。它是似然在先验分布下的加权平均，表示在尚未用当前数据更新参数时，模型赋予这些数据的概率或概率密度。给定模型、先验和数据后，边际似然是一个数；后验则是参数取值的分布。

边际似然首先保证后验归一化：将式\eqref{eq:solving-posterior}对$\theta$积分，分子的积分恰好是$p(\mathcal D)$，所以后验积分为$1$。在本章中，边际似然只作为后验的归一化常数出现。它的模型比较含义以及模型先验的设定，放在后续贝叶斯学习章节讨论。若模型还含有隐变量，计算边际分布时可能需要同时对隐变量和随机参数积分；这一积分通常也没有解析形式。

\subsection{分布近似：从贝叶斯公式推导变分目标}
面对难以解析计算的后验，可以选择一个归一化、便于计算的分布$q_\phi(\theta)$，用它近似$p(\theta\mid\mathcal D)$。例如，可以选择高斯分布，此时$\phi$包含均值和协方差参数。变分推断的任务是确定这些参数，使近似分布尽可能接近后验。下面假设$p(\mathcal D)$正且有限，$q_\phi$在后验为零的区域也为零，且涉及的期望和KL散度有限。

从贝叶斯公式出发，取对数并整理，在后验密度为正的地方有
\begin{equation}
 \log p(\mathcal D)=\log p(\mathcal D,\theta)-\log p(\theta\mid\mathcal D).
\end{equation}
两边对$q_\phi(\theta)$求期望。左边与$\theta$无关，又因为$\int q_\phi(\theta)\,d\theta=1$，所以左边仍为$\log p(\mathcal D)$：
\begin{equation}
 \log p(\mathcal D)=\E_{q_\phi}[\log p(\mathcal D,\theta)]
 -\E_{q_\phi}[\log p(\theta\mid\mathcal D)].
\end{equation}
在右边加上并减去$\E_{q_\phi}[\log q_\phi(\theta)]$，再将各项分组，得到
\begin{align}
 \log p(\mathcal D)
 &=\E_{q_\phi}[\log p(\mathcal D,\theta)-\log q_\phi(\theta)]\notag\\
 &\quad+\E_{q_\phi}[\log q_\phi(\theta)-\log p(\theta\mid\mathcal D)]\notag\\
 &=\mathcal L(\phi)+\KL(q_\phi\|p(\theta\mid\mathcal D)).
 \label{eq:solving-vi-kl}
\end{align}
第一项称为证据下界（Evidence Lower Bound，ELBO）：
\begin{align}
 \mathcal L(\phi)
 &=\E_{q_\phi}[\log p(\mathcal D,\theta)-\log q_\phi(\theta)]\notag\\
 &=\E_{q_\phi}[\log p(\mathcal D\mid\theta)]-\KL(q_\phi\|p(\theta)).
 \label{eq:solving-elbo}
\end{align}
第二项是近似分布与真实后验之间的KL散度。它非负，所以$\mathcal L(\phi)\leq\log p(\mathcal D)$，这解释了“下界”的名称。固定模型时，证据与$\phi$无关，最大化ELBO便等价于最小化这个KL散度。ELBO只涉及联合密度和近似分布，优化时无需计算难以求出的归一化积分。

这里的转化并不声称每个ELBO都能直接算出。计算后验密度需要知道证据$p(\mathcal D)$，而ELBO的第一种写法
\begin{equation}
 \mathcal L(\phi)=\E_{q_\phi}
 \left[\log p(\mathcal D,\theta)-\log q_\phi(\theta)\right]
 \end{equation}
不包含这个未知的归一化常数。只要能够计算给定参数下的联合密度$p(\mathcal D,\theta)=p(\mathcal D\mid\theta)p(\theta)$，并且能够计算或估计关于$q_\phi$的期望，就可以评价或优化ELBO。是否更容易计算，取决于联合密度、近似分布及其期望的具体形式。

若先验KL散度有解析表达式，而对数似然期望难以积分，可以从$q_\phi$抽取样本$\theta^{(1)},\ldots,\theta^{(S)}$，使用
\begin{equation}
 \widehat{\mathcal L}(\phi)
 =\frac1S\sum_{s=1}^S\log p(\mathcal D\mid\theta^{(s)})
 -\KL(q_\phi\|p(\theta))
 \end{equation}
估计ELBO。若先验KL也没有解析形式，则可以使用第一种写法，同时对联合密度和$q_\phi$的对数密度进行蒙特卡洛估计。若某些项可以直接积分或精确计算，应保留解析计算，以减少采样误差。

得到ELBO的数值或无偏（或近似无偏）的梯度估计后，就可以把$\phi$当作普通优化变量。每轮先从当前的$q_{\phi_t}$取样，计算ELBO或其梯度，再按梯度上升更新：
\begin{equation}
 \phi_{t+1}=\phi_t+\alpha_t\widehat{\nabla_\phi\mathcal L(\phi_t)}.
\end{equation}
反复迭代，直到目标变化、参数变化或梯度满足停止条件。求解结果是$q_{\hat\phi}(\theta)$，它是对真实后验的近似，而不是后验的解析表达式。后续的后验均值、区间概率和预测，也都使用这个近似分布计算。

ELBO中的对数似然期望评价近似分布所支持的参数解释数据的程度，先验KL项衡量近似分布与先验的差异。即使优化充分，选定的分布族也可能无法表达真实后验的相关性或多个峰，因此结果仍是近似。这里称为确定性近似，是因为最终用一个确定的分布表达式$q_{\hat\phi}$表示后验；求解其参数时仍可以使用随机采样和随机梯度。

EM算法与变分推断使用了相同的下界恒等式。EM在每次E步计算精确隐变量后验，使变分下界与当前参数下的对数似然相等，再更新固定参数。变分推断将近似分布限制在可处理的族中，下界与对数证据之间通常仍有间隙。这里优化的是分布的参数，求解结果仍是一个分布。

\subsection{均值场近似怎样计算}
均值场（mean-field）近似是变分推断中一种常用的确定性近似。先把参数分成$d$组，记为
\begin{equation}
 \theta=(\theta_1,\ldots,\theta_d),
\end{equation}
并假设各组参数在近似后验中相互独立，于是取变分分布为
\begin{equation}
 q(\theta)=\prod_{j=1}^d q_j(\theta_j).
\end{equation}
其中$q_j(\theta_j)$是只描述第$j$组参数的概率分布，称为一个变分因子；$q_{-j}$表示除$q_j$以外其他因子的乘积。为了求解这些因子，可以固定$q_{-j}$，只把$q_j$作为变量，最大化ELBO关于$q_j$的部分。这里把第$j$个因子看作一个待更新的方向，这种逐个固定其余因子、轮流更新单个因子的做法称为坐标上升变分推断。固定除第$j$个因子以外的$q_{-j}$，令
$A_j(\theta_j)=\E_{q_{-j}}[\log p(\mathcal D,\theta)]$，ELBO中依赖$q_j$的部分为
\begin{equation}
 \int q_j(\theta_j)A_j(\theta_j)\,d\theta_j
 -\int q_j(\theta_j)\log q_j(\theta_j)\,d\theta_j.
\end{equation}
在$\int q_j=1$的约束下，对$q_j$求变分导数，得到$A_j-\log q_j-1+\lambda=0$。归一化以后，坐标更新为
\begin{equation}
 q_j^{\mathrm{new}}(\theta_j)\propto
 \exp\{\E_{q_{-j}}[\log p(\mathcal D,\theta)]\}.
\end{equation}
这要求期望可计算且右端可归一化。逐个更新因子称为坐标上升变分推断。若因子还被限制为指定的参数族，这个不受限的坐标解未必属于该族，可能仍需数值优化。

独立近似降低计算成本，但可能漏掉相关性。例如，真实二维高斯后验协方差为
\begin{equation}
 \Sigma=\begin{pmatrix}1&\rho\\\rho&1\end{pmatrix},\qquad |\rho|<1.
\end{equation}
取独立高斯近似$q=\mathcal N(\bm m,\operatorname{diag}(v_1,v_2))$，最小化$\KL(q\|p)$对均值的解是真实均值。对方差求导得到
$\tfrac12((\Sigma^{-1})_{jj}-1/v_j)=0$，所以$v_1=v_2=1-\rho^2$。真实边际方差均为$1$，而$\rho=0.9$时近似方差只有$0.19$。该例展示了均值场近似可能低估不确定性的原因。增加ELBO表示该目标有所改进，仍需检查分布族能否表达任务所需的相关性、多峰或尾部。

\subsection{蒙特卡洛变分推断与随机梯度优化}
在均值场方法中，如果固定其他因子后，式\eqref{eq:solving-elbo}关于某个因子$q_j$的最优解可以直接写出，就能使用上面的坐标上升更新。复杂模型中，这个条件可能不满足：右端的期望难以计算，或者得到的函数难以归一化；即使能得到完整的最优分布，它也可能不属于预先选定的高斯等分布族。此时，不能直接写出$q_j$的更新式，通常改为把近似分布写成$q_\phi(\theta)$，再用数值优化确定其参数。此时ELBO中的期望也可能难以直接积分，可以从当前近似分布$q_\phi$抽取样本，用样本平均估计ELBO或其梯度。这种做法称为蒙特卡洛变分推断（Monte Carlo Variational Inference，MCVI）。待求的仍是近似分布的参数$\phi$，采样为优化提供可计算的估计。

令$f_\phi(\theta)=\log p(\mathcal D,\theta)-\log q_\phi(\theta)$，则$\mathcal L(\phi)=\E_{q_\phi}[f_\phi(\theta)]$。固定当前$\phi$，独立抽取$\theta^{(s)}\sim q_\phi$，可用$S^{-1}\sum_{s=1}^S f_\phi(\theta^{(s)})$估计ELBO。在可积条件下，这个估计的期望等于真实ELBO。不过，$q_\phi$本身随$\phi$变化，计算梯度时也需要考虑采样分布的变化。如果把抽出的$\theta^{(s)}$直接视为与$\phi$无关的常数再求导，就会遗漏这部分影响。下面介绍两种常见的梯度估计方式。

\paragraph{重参数化梯度}
对角高斯近似可以写成
\begin{equation}
 \theta=\bm\mu_\phi+\bm s_\phi\odot\epsilon,
 \qquad \epsilon\sim\mathcal N(\bm0,I),
 \label{eq:solving-reparameterization}
\end{equation}
其中$\bm s_\phi$逐坐标为正，可用对数尺度参数保证正值。随机性被放入与$\phi$无关的$\epsilon$，称为重参数化。抽取$S$个独立样本，令$\theta^{(s)}=\bm\mu_\phi+\bm s_\phi\odot\epsilon^{(s)}$，可估计
\begin{equation}
 \widehat{\mathcal L}(\phi)=\frac1S\sum_{s=1}^S
 [\log p(\mathcal D,\theta^{(s)})-\log q_\phi(\theta^{(s)})].
\end{equation}
在可微且允许交换求导与期望的条件下，对这个重参数化后的表达式自动微分，可以得到ELBO梯度的蒙特卡洛估计。随后进行梯度上升，或最小化负ELBO。高斯熵或先验KL若可解析计算，应直接计算这些项，以减少采样噪声。

\paragraph{得分函数梯度}
如果随机变量是离散的，或者难以写出重参数化表达式，可以使用恒等式
\begin{equation}
 \nabla_\phi\mathbb E_{q_\phi}[f(\theta)]
 =\mathbb E_{q_\phi}\left[f(\theta)\nabla_\phi\log q_\phi(\theta)
 +\nabla_\phi f(\theta)\right].
 \label{eq:solving-score-function-gradient}
\end{equation}
这个恒等式要求分布的支持集不随$\phi$变化，并允许交换求导与积分（离散变量时为求和）。其中$\nabla_\phi\log q_\phi(\theta)$称为得分函数。对于ELBO中的$f_\phi$，固定生成模型时有$\nabla_\phi f_\phi=-\nabla_\phi\log q_\phi$，而得分函数的期望为零，因此
\begin{equation}
 \nabla_\phi\mathcal L(\phi)
 =\E_{q_\phi}[f_\phi(\theta)\nabla_\phi\log q_\phi(\theta)].
\end{equation}
用独立样本替换期望，可得梯度估计
\begin{equation}
 \widehat{\bm g}(\phi)=\frac1S\sum_{s=1}^S
 [f_\phi(\theta^{(s)})-b]\nabla_\phi\log q_\phi(\theta^{(s)}).
\end{equation}
这里$b$是与本次样本无关的基线，例如由以往迭代估计的常数。减去$b$不改变梯度估计的期望，因为$\E_{q_\phi}[b\nabla_\phi\log q_\phi]=0$，但合适的基线可以降低方差。在上述条件下，这个估计是无偏的。得分函数方法适用范围较广，方差却可能较大；重参数化梯度在许多可微模型中方差较小，但需要合适的可微采样表达式。

\paragraph{用随机梯度更新变分参数}
将蒙特卡洛梯度估计用于优化，便得到随机梯度变分推断。第$t$次迭代从当前$q_{\phi_t}$重新采样，计算梯度估计$\widehat{\bm g}_t$，再按步长$\alpha_t$更新：
\begin{equation}
 \phi_{t+1}=\phi_t+\alpha_t\widehat{\bm g}_t.
\end{equation}
完整流程是选择分布族、初始化变分参数、采样计算ELBO及梯度、更新变分参数，再检查目标变化与数值稳定性。MCVI描述用蒙特卡洛方法计算变分目标或梯度，随机梯度变分推断描述利用这些梯度估计进行优化。即使使用全部训练数据，从$q_\phi$采样也会使梯度具有随机性。

如果给定参数后各样本条件独立，大数据下还可以均匀随机抽取小批量$B_t$，用下面的量估计全数据的对数似然总和：
\begin{equation}
 \widehat\ell_{B_t}(\theta)
 =\frac{n}{|B_t|}\sum_{i\in B_t}\log p(y_i\mid\bm x_i,\theta).
\end{equation}
这里缩放的是数据项，先验项和近似分布的熵项保持原有权重。参数采样和数据抽样是两种不同的随机来源。采样目标有噪声，单次估计增加不表示真实ELBO必然增加，优化收敛也仍依赖步长、梯度估计的性质和目标函数。

离散变量通常没有上述直接可微的重参数化，应先检查能否精确求和，或采用适合模型结构的其他梯度估计。最终预测可以从$q_\phi$抽取参数，平均各次条件预测。这个平均仍包含后验近似误差与有限采样误差。

若模型还包含确定性参数$\eta$，例如表示条件分布的神经网络权重，则联合密度写为$p_\eta(\mathcal D,\theta)$，变分目标也随之写为$\mathcal L(\phi,\eta)$。可以交替更新近似后验的参数$\phi$和神经网络参数$\eta$，也可以用梯度方法联合更新。在近似后验受限时，优化的是对数边际似然的下界；只有在近似分布等于对应的真实后验时，下界才与对数边际似然相等。

\subsection{随机性近似：用后验样本计算期望}
另一个求解角度是从最终需要计算的量出发。若关注参数均值、区间概率或预测，通常需要的是某个函数$h(\theta)$关于后验的期望：
\begin{equation}
 \E[h(\theta)\mid\mathcal D]
 =\int h(\theta)p(\theta\mid\mathcal D)\,d\theta.
\end{equation}
例如，$h(\theta)=\theta$给出后验均值，$h(\theta)=\mathbf 1_{\{\theta\in A\}}$给出参数落在区域$A$内的后验概率。将$h(\theta)$取为条件预测概率或密度，还可以计算式\eqref{eq:solving-posterior-predictive}中的后验预测。以下要求所求期望存在。

如果能够从后验分布抽取独立样本$\theta^{(1)},\ldots,\theta^{(S)}$，就可以用蒙特卡洛平均近似积分：
\begin{equation}
 \E[h(\theta)\mid\mathcal D]
 \approx\frac1S\sum_{s=1}^S h(\theta^{(s)}).
\end{equation}
大数定律说明，在可积条件下，随着样本数增加，这个平均趋于目标期望。这样可以用同一批参数样本计算不同函数的期望，无需先写出后验的解析表达式。若样本来自变分分布$q_\phi$，平均对应的是$q_\phi$下的期望，仍会保留变分近似误差。

难点在于，无法解析计算的后验通常也不能直接进行独立采样。马尔可夫链蒙特卡洛（Markov Chain Monte Carlo，MCMC）通过构造以后验为平稳分布的马尔可夫链来解决这一问题。以Metropolis--Hastings算法为例，令未归一化密度$\widetilde p(\theta)=p(\mathcal D\mid\theta)p(\theta)$。从当前状态$\theta$按提议密度$r(\theta'\mid\theta)$产生候选状态$\theta'$，以概率
\begin{equation}
 \alpha(\theta,\theta')=\min\left\{1,
 \frac{\widetilde p(\theta')r(\theta\mid\theta')}
 {\widetilde p(\theta)r(\theta'\mid\theta)}\right\}
\end{equation}
接受候选状态，否则保留当前状态。后验的归一化常数在比值中相消，因此不必先计算$p(\mathcal D)$。这一接受规则使后验满足详细平衡，从而成为链的平稳分布。在适当的不可约性、非周期性等条件下，链能够趋于目标后验，遍历平均可以近似后验期望。

链的样本通常相关，有限运行也可能尚未充分探索后验。实际计算需要检查初始阶段的影响、多条链之间的一致性及有效样本量。增加迭代次数只有在链能够充分探索目标分布时才有意义，不能把迭代数直接当作独立样本数。后验采样用于估计后验期望，也不意味着已经求出了模型比较所需的边际似然。

这里的后验采样称为蒙特卡洛推断，其中MCMC是一类常用方法。蒙特卡洛变分推断（Monte Carlo Variational Inference，MCVI）则用采样估计变分目标或其梯度，仍然是在优化近似分布$q_\phi$。两者使用采样的目的不同：前者用以后验为目标的样本近似期望，后者用采样辅助求解变分近似。后续贝叶斯学习章节将进一步讨论这些推断方法。

\section{选择求解方法时应问的问题}

模型求解前，应确认待确定的是固定参数、隐变量的条件分布，还是参数后验。然后检查线性方程、期望或积分能否直接计算，以及目标能否求导。

固定参数有可靠直接解时，可以使用数值线性代数。无直接解且可微时，可以使用梯度方法。无法取得梯度的离散或黑盒问题，可以在计算预算允许时采用启发式搜索。含隐变量的极大似然问题在条件期望可处理时适合EM。随机参数的后验不可解析时，变分推断将分布近似转化为优化问题。

这些方法的结果具有不同含义。目标下降评价参数搜索，似然提高评价数据在模型下的概率，ELBO提高评价所选分布族中的近似目标。它们都不能替代下一章的泛化检验。

\section*{辅助学习材料}
\begin{itemize}
\item \textbf{Mathematics for Machine Learning}；Deisenroth、Faisal与Ong；英文教材。结合线性代数、微积分和连续优化部分阅读最小二乘与梯度方法。\cite{deisenroth2020mml}
\item \textbf{Probabilistic Machine Learning: An Introduction}；Murphy；英文教材。对照混合模型与潜变量推断，理解条件分布在模型求解中的作用。\cite{murphy2022pml}
\item \textbf{Pattern Recognition and Machine Learning}；Bishop；英文教材。第九章介绍混合模型与EM，第十章介绍变分推断。\cite{bishop2006pattern}
\end{itemize}

\section{本章小结}

模型求解首先需要明确未知量。固定参数可以通过闭式计算、梯度方法或启发式搜索获得点估计。含隐变量的极大似然估计利用EM交替计算条件期望和更新参数。随机参数模型需要后验分布，变分推断通过优化ELBO得到可计算的近似。

各方法的保证都有条件。梯度下降的单步保证依赖光滑性与步长，EM的单调性依赖E步和M步的计算，变分近似质量受分布族限制。后续统计学习内容将讨论怎样用独立数据检验所得模型能否用于未见样本。



